\documentclass[10pt,a4paper]{amsart}
\usepackage[margin=19mm]{geometry}
\usepackage{amsmath,amssymb,mathtools}
\usepackage[scr=rsfs,cal=euler]{mathalfa}
\usepackage{xfrac}
\usepackage[unicode,hidelinks,bookmarks=false]{hyperref}
\newcommand{\ie}{i.e.\ }
\newcommand{\cf}{cf.\ }
\newcommand{\resp}{respectively\ }
\newcommand{\Subsection}{\subsection}
\providecommand{\nobreakdash}{\nobreak\hbox{-}}
\setlength{\emergencystretch}{2em}
\multlinegap0pt
\usepackage{amssymb}
\usepackage{esint}
\usepackage{bm}
\usepackage{xcolor}
\newtheorem{remark}{Remark}
\newtheorem{lemma}{Lemma}
\title{Global Weighted Gradient Estimates for Double Obstacle Problems with Orlicz Growth}
\author{Qi Xiong, Xing Fu}
\date{Source: arXiv:2608.27935v1 (CC BY 4.0)}
\begin{document}
\maketitle

\noindent\emph{Source and licence.} Global Weighted Gradient Estimates for Double Obstacle Problems with Orlicz Growth. Version arXiv:2608.27935v1 (CC BY 4.0), distributed by arXiv under the Creative Commons Attribution 4.0 International licence, http://creativecommons.org/licenses/by/4.0/, as recorded in the arXiv licence metadata for this exact version and shown on the abstract page https://arxiv.org/abs/2608.27935v1. Full licence notice: \texttt{licenses/arxiv-2608.27935-CC-BY-4.0.txt}.\par\smallskip

\noindent\textit{Editorial scope.} Counted unit: the article's lemma lem5.1 (source0.tex lines 905--925). The article's complete original proof of it is delivered with it (source0.tex lines 927--1159, one \texttt{proof} environment of 233 lines). No mathematics is written, repaired or re-derived: the statement and the proof are the article's own lines, in the article's own notation, and no retained line is substituted or re-flowed.  Retained uncounted as the source-local context the proof needs: lines 176--203; lines 270--279; lines 393--395; lines 472--476; lines 538--540; lines 542--544; lines 549--551; lines 553--569; lines 689--704; lines 807--821; lines 860--874; lines 882--896.  Nothing was omitted from any retained span, so the package carries no omission record.  No retained line contains a citation, so the wrapper carries no bibliography and no bibliography entry.  No retained line contains a figure environment, an image-inclusion command or drawing code, so no image is delivered and the package declares no asset dependency.  Editorial formatting only, in addition: an AMS \texttt{amsart} preamble that restates the article's own declarations and macros used by the retained lines, namely the environments \texttt{remark}, \texttt{lemma} on separate counters, together with the standard packages those lines need.  The retained environments print in this document as Remark 1, Lemma 1 in order of appearance, whereas the article prints the same items with the numbering of its own full document.  The complete native source of the article as distributed in the arXiv e-print for this exact version is bundled unchanged as \texttt{sources/208748/source0.tex}.
\begin{equation}\label{u}
	\int_{\Omega}  a(x, D u) \cdot D(v-u) \mathrm{d} x \geqslant \int_{\Omega} b(x, F ) \cdot D(v-u) \mathrm{d} x  
\end{equation}
for all $v \in \mathcal{A}_{h}(\Omega)$. Here we
assume that $F \in L^G(\Omega;\mathbb{R}^n)$ and  $a=a(x,\eta): \Omega \times\mathbb{R}^n \rightarrow \mathbb{R}^n$ is measurable for every $\eta \in \mathbb{R}^n$  and differentiable for almost every $x\in \Omega$  and there exist constants $0<l\leqslant 1 \leqslant L<+\infty$ such that for all $x \in \Omega,\eta,\lambda \in \mathbb{R}^n$,   
\begin{eqnarray} \label{u1}
	\left\{\begin{array}{r@{}c@{}ll}
		&&D_{\eta} a(x,\eta )\lambda \cdot \lambda \geqslant l\dfrac{g(|\eta|)}{|\eta|}|\lambda|^2 \,, \\[0.05cm]
		&&|a(x,\eta)|+|\eta||D_{\eta} a(x,\eta )|\leqslant Lg(|\eta|)\,, \\[0.05cm]
	\end{array}\right.
\end{eqnarray}
where $D_{\eta}$ denotes the differentiation in $\eta$.  
Moreover, let $b(x,\eta): \Omega \times\mathbb{R}^n \rightarrow \mathbb{R}^n$  satisfy
\begin{equation}\label{u2}
	|b(x,\eta)|\leqslant L g(|\eta|).
\end{equation}
Here $g(t) : [0,+\infty)\rightarrow [0,+\infty)$ satisfies
\begin{eqnarray}\label{u3}
	\left\{\begin{array}{r@{}c@{}ll}
		&&g(t)=0 \ \ \ \Leftrightarrow \ \ \   t=0 \,, \\[0.05cm]
		&&g(\cdot)\in C^{1}(\mathbb{R}^+)\,, \\[0.05cm]
		&&   i_{g}=: \inf_{t>0}\frac{tg'(t)}{g(t)}\leq \sup_{t>0}\frac{tg'(t)}{g(t)}=:s_{g}<\infty, \ \ \  where \ \  0< i_{g} \leqslant 1 \leqslant s_{g} <\infty. \, \\[0.05cm]
	\end{array}\right.
\end{eqnarray}
Define
\begin{equation}\label{g}
	G(t):= \int_0^tg(\tau)\,d\tau  \ \ \ \mbox{for}\ \ t\geq0.
\end{equation}

\begin{remark}\label{rei}
	For every $x \in \partial \Omega$ and $r \in (0,R]$ we choose a coordinate system $(z_1,\dots,z_n)$ with the origin at some interior point of $  \Omega$ such that in this coordinate system $x=- \delta re_n$ and
	$$
	B_{\frac{7}{8}r}(0)\cap \{z_n>0\} \ \subset\ \Omega\cap B_{\frac{7}{8}r}(0)\ \subset\ B_{\frac{7}{8}r}(0)\cap \{z_n>-2\delta r\}.
	$$
	
	
	Note that a domain with Lipschitz boundary with Lipschitz seminorm $\delta\in(0,1)$
	is $(\delta,R)$-Reifenberg flat for some $R>0$ and that the ball $B_r$ is $(\delta,2\delta r)$-Reifenberg flat for any $\delta\in (0,\frac{1}{2})$.
\end{remark}

 \begin{align}\label{cv}
G(|\xi-\eta|)\leqslant c_v \left[ |V(\xi) - V(\eta)|^2+G(\eta)\right] 
	\end{align}

\begin{align}\label{one}
	\left|\left\{ x\in\mathbb{R}^n : \mathcal{M}f(x) > K \right\}\right|
	&\leq \frac{c}{K} \int_{\mathbb{R}^n}|f(x)|\,\mathrm{d}x,
	\qquad \text{for every } K>0,
\end{align}

\begin{equation}\label{ass1}
  \sup_{0<\varrho<5} \fint_{B_\varrho} \theta(a, B_\varrho)(x) \, \mathrm{d}x \leq \delta,
\end{equation}

\begin{equation}\label{ass2}
\fint_{\Omega_5} G(|Du|) \, \mathrm{d}x \leq 1, \qquad \fint_{\Omega_5}\left[  G(|F|)+G(|Dh|)+G(|D\psi_1|)+G(|D\psi_2|)\right]  \, \mathrm{d}x \leq \delta,
\end{equation}

\begin{align}\label{ass3}
 B_5^+ \subset \Omega_5 \subset \left\{ x \in B_5 : x_n > -\frac{80}{7}\delta \right\},
\end{align}

\begin{lemma}\label{uw1}
	 Assume that \eqref{u1}-\eqref{u3} are satisfied and   $u \in \mathcal{A}_{h}(\Omega)$   satisfies the variational inequality \eqref{u}.
	 For any $\varepsilon > 0$ there exists a small $\delta=\delta(n,l,L,i_g,s_g,\varepsilon) \in (0,\frac{1}{8})$ such that if \eqref{ass2} holds for such $\delta$ and  $w_1 \in W^{1,G}(\Omega_5)$ with $\psi_1-h\leqslant w_1\leqslant \psi_2-h $ a.e. in $\Omega_5$ solves
	 \begin{equation}\label{w1}
	 	\left\{\begin{array}{r@{\ \ }c@{\ \ }ll}
	 		&\int_{\Omega_5}  a(x,  Dw_1) \cdot D(v-w_1)dx \geqslant \int_{\Omega_5}  b(x, F) \cdot D(v-w_1)dx,  \\[0.05cm]
	 		&w_1=u-h \ \ \mbox{on}\ \ \partial \Omega_5 \,,
	 	\end{array}\right.
	 \end{equation}
	where $v\in w_1+W_0^{1,G}(\Omega_5)$ satisfies $\psi_1-h \leqslant v \leqslant \psi_2-h$ a.e. in $\Omega_5$.
	Then
	\begin{equation}\label{vuw1}
		\fint_{\Omega_5} G(|Dw_1|)\,dx \leq c \quad \text{and} \quad \fint_{\Omega_5} |V(Du)-V(Dw_1)|^2\,dx \leq \varepsilon.  
	\end{equation}
	Here, $c > 0$ depends on $n,l,L,i_g$ and $s_g$, but is independent of $\varepsilon$.
	
\end{lemma}

\begin{lemma}\label{w1w2l}
	Assume that \eqref{u1}-\eqref{u3} are satisfied and   $w_1 \in u-h+ W_0^{1,G}(\Omega_5)$ with $\psi_1-h\leqslant w_1\leqslant \psi_2-h $ a.e. in $\Omega_5$ solves the inequality \eqref{w1}.  
	For any $\varepsilon > 0$ there exists a small $\delta=\delta(n,l,L,i_g,s_g,\varepsilon) \in (0,\frac{1}{8})$ such that if \eqref{ass2} holds for such $\delta$ and  $w_2 \in W^{1,G}(\Omega_5)$ with $ w_2 \geqslant \psi_1-h$ a.e. in $\Omega_5$ solves
	\begin{equation}\label{w2}
		\left\{\begin{array}{r@{\ \ }c@{\ \ }ll}
			&\int_{\Omega_5}  a(x,  Dw_2) \cdot D(v-w_2)dx \geqslant \int_{\Omega_5}  a(x,  D\psi_2-Dh) \cdot D(v-w_2)dx,  \\[0.05cm]
			&w_2=w_1 \ \ \mbox{on}\ \ \partial \Omega_5 \,,
		\end{array}\right.
	\end{equation}
	where $v\in w_2+W_0^{1,G}(\Omega_5)$ satisfies $ v \geqslant \psi_1-h$ a.e. in $\Omega_5$.
	Then
	\begin{equation}\label{w1w2}
		\fint_{\Omega_5} G(|Dw_2|)\,dx \leq c \quad \text{and} \quad \fint_{\Omega_5} |V(Dw_1)-V(Dw_2)|^2\,dx \leq \varepsilon.  
	\end{equation}
	Here, $c > 0$ depends on $n,l,L,i_g$ and $s_g$, but is independent of $\varepsilon$.
	\end{lemma}

		\begin{lemma}\label{w2w3l}
		Assume that \eqref{u1}-\eqref{u3} are satisfied and   $w_2 \in w_1+ W_0^{1,G}(\Omega_5)$ with $ w_2 \geqslant \psi_1-h$ a.e. in $\Omega_5$ solves  the inequality \eqref{w2}.  
		For any $\varepsilon > 0$ there exists a small $\delta=\delta(n,l,L,i_g,s_g,\varepsilon) \in (0,\frac{1}{8})$ such that if \eqref{ass2} holds for such $\delta$ and  $w_3 \in W^{1,G}(\Omega_5)$ is a weak solution of 
		\begin{equation}\label{w3}
			\left\{\begin{array}{r@{\ \ }c@{\ \ }ll}
					-\operatorname{div}a(x,Dw_3) &=&-\operatorname{div}a(x,D\psi_1-Dh) \ \ &\mbox{in}\ \   \Omega_5 \,,  \\[0.05cm]
				 w_3 &=&w_2 \ \ &\mbox{on}\ \ \partial \Omega_5 \,.
			\end{array}\right.
		\end{equation}
		Then
		\begin{equation}\label{w2w3}
			\fint_{\Omega_5} G(|Dw_3|)\,dx \leq c \quad \text{and} \quad \fint_{\Omega_5} |V(Dw_2)-V(Dw_3)|^2\,dx \leq \varepsilon.  
		\end{equation}
		Here, $c > 0$ depends on $n,l,L,i_g$ and $s_g$, but is independent of $\varepsilon$.
	\end{lemma}

	\begin{lemma}\label{w3w4l}
	Assume that \eqref{u1}-\eqref{u3} are satisfied and   $w_3 \in w_2+ W_0^{1,G}(\Omega_5)$  is a weak solution of \eqref{w3}. 
	For any $\varepsilon > 0$ there exists a small $\delta=\delta(n,l,L,i_g,s_g,\varepsilon) \in (0,\frac{1}{8})$ such that if \eqref{ass1}  and \eqref{ass2}   holds for such $\delta$ and  $w_4 \in W^{1,G}(B_4)$ is a weak solution of 
	\begin{equation*} 
		\left\{\begin{array}{r@{\ \ }c@{\ \ }ll}
			-\operatorname{div}\bar{a}_{B_4}(Dw_4) &=&0 \ \ &\mbox{in}\ \   B_4 \,,  \\[0.05cm]
			w_4 &=&w_3 \ \ &\mbox{on}\ \ \partial B_4 \,.
		\end{array}\right.
	\end{equation*}
	Then
	\begin{equation*} 
		\|G(|Dw_4|)\|_{L^{\infty}(B_2)} \leq c \quad \text{and} \quad \fint_{B_4} |V(Dw_3)-V(Dw_4)|^2\,dx \leq \varepsilon.  
	\end{equation*}
	Here, $c > 0$ depends on $n,l,L,i_g$ and $s_g$, but is independent of $\varepsilon$.
\end{lemma}

		\begin{lemma}\label{w3w4b}
		Assume that \eqref{u1}-\eqref{u3} are satisfied and   $w_3 \in w_2+ W_0^{1,G}(\Omega_5)$  is a weak solution of \eqref{w3}. 
		For any $\varepsilon > 0$ there exists a small $\delta=\delta(n,l,L,i_g,s_g,\varepsilon) \in (0,\frac{1}{8})$ such that if \eqref{ass1},  \eqref{ass2} and \eqref{ass3} holds for such $\delta$ and  $w_4 \in W^{1,G}(B_4^+)$ is a weak solution of 
		\begin{equation*} 
			\left\{\begin{array}{r@{\ \ }c@{\ \ }ll}
				-\operatorname{div}\bar{a}_{B_4^+}(Dw_4) &=&0 \ \ &\mbox{in}\ \   B_4^+ \,,  \\[0.05cm]
				w_4 &=&0 \ \ &\mbox{on}\ \   T_4 \,.
			\end{array}\right.
		\end{equation*}
		Then
		\begin{equation*} 
			\|G(|D\bar{w}_4|)\|_{L^{\infty}(\Omega_2)} \leq c \quad \text{and} \quad \fint_{\Omega_3} |V(Dw_3)-V(D\bar{w}_4)|^2\,dx \leq \varepsilon,  
		\end{equation*}
	where $\bar{w}_4$ is the zero extension of $w_4$ from $B_4^+$ to $B_4$ and $c > 0$ depends on $n,l,L,i_g$ and $s_g$, but is independent of $\varepsilon$.
	\end{lemma}

 \begin{lemma}\label{lem5.1}
 	There exists a constant $C = C(n,l,L,i_g, s_g ) > 1$ such that for any $\varepsilon > 0$, there exists $\delta = \delta(\varepsilon,n,l,L,i_g, s_g) \in (0, \frac{1}{8})$ ensuring that if $(a, \Omega)$ satisfies the $(\delta, R)$-vanishing condition, then the following holds. For any $y \in \Omega$, $r \in (0, \frac{R}{40})$, and $\lambda \geq 1$, provided that the set 
 	\begin{align}\label{ml}
 		&\left\{
 		\begin{aligned}
 			x\in\Omega_r(y):\quad
 			&\mathcal{M}\bigl(G(|Du|)\bigr)(x)\leq\lambda,\quad
 			\mathcal{M}\bigl(G(|F|)\bigr)(x)\leq\frac{\lambda\delta}{4},\quad
 			\mathcal{M}\bigl(G(|D\psi_1|)\bigr)(x)\leq\frac{\lambda\delta}{4},
 			\\
 			&\mathcal{M}\bigl(G(|D\psi_2|)\bigr)(x)\leq\frac{\lambda\delta}{4},\quad
 			\mathcal{M}\bigl(G(|Dh|)\bigr)(x)\leq\frac{\lambda\delta}{4}
 		\end{aligned}
 		\right\}
 		\neq\varnothing,
 	\end{align}
 	then
 	\begin{align}
 		|\{x \in \Omega_r(y) : \mathcal{M}(G(|Du|)) > \lambda C\}| < \varepsilon |B_r(y)|.
 	\end{align}
 \end{lemma}

 \begin{proof}
 	By our assumption in \eqref{ml}, there exists a point $y_1 \in \Omega_r(y)$ such that for every $\rho > 0$, we have
  \begin{equation}\label{ml2}
  	\left\{
  	\begin{aligned}
  		\frac{1}{|B_\rho(y_1)|}
  		\int_{\Omega_\rho(y_1)}G(|Du|)\,dx
  		&\leq \lambda, \\[0.15cm]
  		\frac{1}{|B_\rho(y_1)|}
  		\int_{\Omega_\rho(y_1)}G(|F|)\,dx
  		&\leq \frac{\lambda\delta}{4}, \\[0.15cm]
  		\frac{1}{|B_\rho(y_1)|}
  		\int_{\Omega_\rho(y_1)}G(|D\psi_i|)\,dx
  		&\leq \frac{\lambda\delta}{4},	\qquad i=1,2,  \\[0.15cm]
  		\frac{1}{|B_\rho(y_1)|}
  		\int_{\Omega_\rho(y_1)}G(|Dh|)\,dx
  		&\leq \frac{\lambda\delta}{4}.
  	\end{aligned}
  	\right.
  \end{equation}
 	Let us first consider the case where $B_{5r}(y) \subset \Omega$. This geometric condition ensures the inclusion $B_{5r}(y) \subset \Omega_{6r}(y_1)$. Combining this fact with the estimates in \eqref{ml2} yields
 	\begin{equation}\label{ml3}
 		\left\{
 		\begin{aligned}
 				\fint_{B_{5r}(y)} G(|Du|) \, dx  &\leq 2^n \lambda,  \\[0.15cm]
 		 \fint_{B_{5r}(y)} G(|F|) \, dx &\leq \frac{ 2^n \lambda \delta}{4}, \\[0.15cm]
 		  \fint_{B_{5r}(y)} G(|D\psi_i|) \, dx &\leq \frac{ 2^n \lambda \delta}{4},	\qquad i=1,2,  \\[0.15cm]
 		\fint_{B_{5r}(y)} G(|Dh|) \, dx &\leq \frac{ 2^n \lambda \delta}{4}.
 		\end{aligned}
 		\right.
 	\end{equation}
 	To proceed, we introduce the following rescaling argument.
 		\begin{equation*} 
 		\left\{
 		\begin{aligned}
 			 	&\tilde{u}(x) = \frac{u(y + rx)}{2^n \lambda r}, \quad \tilde{F}(x) = \frac{F(y + rx)}{2^n \lambda},  \quad \tilde{\psi_i}(x) = \frac{\psi_i(y + rx)}{2^n \lambda r}, \quad  i=1,2, \\[0.15cm]
 			 	&\tilde{a}(x,\eta) = a(y + rx,2^n \lambda \eta), \quad 	\tilde{b}(x,\eta) = b(y + rx,2^n \lambda \eta), \quad \tilde{h}(x) = \frac{h(y + rx)}{2^n \lambda r},   \\[0.15cm]
 		&\tilde{g}(t) = g(2^n \lambda t), \quad \tilde{G}(t) =\int_0^t \tilde{g}(s) ds, \quad t \geq 0, \\[0.15cm]
 		\end{aligned}
 		\right.
 	\end{equation*}
 	where
 	\[
 	x \in \tilde{\Omega} = \left\{ \frac{z - y}{r} : z \in \Omega \right\}. 
 	\]
 	Then $\tilde{a}, \tilde{b}$ and $\tilde{g}$ satisfy \eqref{u1}, \eqref{u2}, \eqref{u3}, and $\tilde{u}\in \tilde{\mathcal{A}}_{\tilde{h}}(\tilde{\Omega})$ solves the variational inequality
 	\begin{equation*} 
 		\int_{\tilde{\Omega}}  \tilde{a}(x, D \tilde{u}) \cdot D(v-\tilde{u}) \mathrm{d} x \geqslant \int_{\tilde{\Omega}} \tilde{b}(x, \tilde{F} ) \cdot D(v-\tilde{u}) \mathrm{d}x,  
 	\end{equation*}
 	where 
 	$$
 	\tilde{\mathcal{A}}_{\tilde{h}}(\tilde{\Omega})=\left\{\varphi \in W_{\tilde{h}}^{1, \tilde{G}}(\tilde{\Omega}): \tilde{\psi}_{1} \leq \varphi \leq \tilde{\psi}_{2} \text { a.e. in } \tilde{\Omega} \right\} .
 	$$
 	   Therefore, conditions \eqref{ass1} and \eqref{ass2} are satisfied, and $\frac{R}{r}>5$. By applying Lemmas \ref{uw1}, \ref{w1w2l}, \ref{w2w3l}, and \ref{w3w4l} and scaling back, we obtain a function $w_4 \in W^{1,G}(B_{4r}(y))$ such that
 	\begin{align}\label{ml4}
 			\fint_{B_{4r}(y)} |V(Du) - V(Dw_4)|^2 \, dx < 2^n \lambda \varepsilon_1 \quad \text{and} \quad \|G(|Dw_4|)\|_{L^\infty(B_{2r}(y))} \leq 2^n \lambda C_1, 
 	\end{align}
 	where $\varepsilon_1>0$ is to be chosen later, and $C_1>1$ is a constant depending only on $n,l,L,i_g,s_g$.
 	Set
 	\[
 	C\geq \max\{2^{n+1}c_vC_1,3^n\}
 	\]
 	for $c_v$ as in \eqref{cv}. We next claim the following inclusion:
 	\begin{align}\label{da}
 		\left\{x\in B_r(y):\mathcal{M}(G(|Du|))>\lambda C\right\}
 		\subset
 		\left\{x\in B_r(y):
 		\mathcal{M}_{B_{4r}(y)}
 		\left(|V(Du)-V(Dw_4)|^2\right)>2^n\lambda
 		\right\}.
 	\end{align}
 To this end, let us take a point
 	\[
 	x_1\in
 	\left\{
 	x\in B_r(y):
 	\mathcal{M}_{B_{4r}(y)}
 	\left(|V(Du)-V(Dw_4)|^2\right)
 	\leq 2^n\lambda
 	\right\}.
 	\]
 First, we consider the case where $0<\rho<r$. Then $B_\rho(x_1)\subset B_{2r}(y)$, and combining \eqref{cv} with \eqref{ml4} yields
 	\begin{align*}
 		\frac{1}{|B_\rho(x_1)|}
 		\int_{\Omega_\rho(x_1)}G(|Du|)\,dx
 		&\leq
 		c_v
 		\fint_{B_\rho(x_1)}
 		\left(
 		|V(Du)-V(Dw_4)|^2
 		+
 		G(|Dw_4|)
 		\right)dx
 		\nonumber\\
 		&\leq
 		c_v(2^n\lambda+2^n\lambda C_1)
 		\leq \lambda C.
 	\end{align*}
  Next, we consider the case where $\rho\geq r$. In this situation, the geometric inclusion $B_\rho(x_1)\subset B_{3\rho}(y_1)$ combined with \eqref{ml2} implies that
 	\begin{align*}
 		\frac{1}{|B_\rho(x_1)|}
 		\int_{\Omega_\rho(x_1)}G(|Du|)\,dx
 		&\leq
 		\frac{|B_{3\rho}(y_1)|}{|B_\rho(x_1)|}
 		\frac{1}{|B_{3\rho}(y_1)|}
 		\int_{\Omega_{3\rho}(y_1)}G(|Du|)\,dx
 		\nonumber\\
 		&\leq 3^n\lambda
 		\leq \lambda C .
 	\end{align*}
 	 Consequently, we obtain
 	\begin{align*}
 		x_1\in
 		\left\{
 		x\in B_r(y):
 		\mathcal{M}(G(|Du|))\leq \lambda C
 		\right\},
 	\end{align*}
which establishes the claimed inclusion \eqref{da}.
 	
 Now, by virtue of \eqref{one}, \eqref{ml4}, and \eqref{da}, we can estimate
 	\begin{align*}
 		\left|
 		\left\{
 		x\in B_r(y):
 		\mathcal{M}(G(|Du|))>\lambda C
 		\right\}
 		\right|
 		&\leq
 		\left|
 		\left\{
 		x\in B_r(y):
 		\mathcal{M}_{B_{3r}(y)}
 		(|V(Du)-V(Dw_4)|^2)>2^n\lambda
 		\right\}
 		\right|
 		\nonumber\\
 		&\leq
 		\frac{c|B_r(y)|}{2^n\lambda}
 		\fint_{B_{3r}(y)}
 		|V(Du)-V(Dw_4)|^2\,dx
 		\nonumber\\
 		&<
 		c\varepsilon_1|B_r(y)|.
 	\end{align*}
 Finally, the proof is completed by choosing $\varepsilon_1$ sufficiently small such that $c\varepsilon_1\leq \varepsilon$.
 	
 	
 	Next, we address the boundary case where $B_{5r}(y) \nsubseteq \Omega$, meaning there exists a point $\bar{x} \in B_{5r}(y) \cap \partial\Omega$. In light of Remark \ref{rei} and the assumption $0 < 40r < R$, we can select a local coordinate system $\{\xi_1, \ldots, \xi_n\}$ such that $\bar{x} = -40r\delta e_n$ and the following geometric inclusion holds:
 	\begin{align}\label{rei2}
 		B^+_{35r}\subset \Omega_{35r}
 		\subset
 		\{z\in B_{35r}:z_n>-80r\delta\}.
 	\end{align}
 	Within this coordinate system, we have $|y| \leq 10r$ and $|y_1| \leq 11r$. This implies the inclusions $\Omega_r(y) \subset \Omega_{11r}$ and $\Omega_{35r} \subset \Omega_{46r}(y_1)$. Consequently, it follows from \eqref{ml2} that
 		\begin{equation}\label{rei3}
 		\left\{
 		\begin{aligned}
 			 	\fint_{\Omega_{35r}}G(|Du|)\,dz
 			 &\leq 2^{n+1}\lambda, \\[0.15cm]
 		 \fint_{\Omega_{35r}}G(|F|)\,dz
 		 &\leq \frac{2^{n+1}\lambda\delta}{4}, \\[0.15cm]
 		 \fint_{\Omega_{35r}}G(|Dh|)\,dz
 		 &\leq \frac{2^{n+1}\lambda\delta}{4}, \\[0.15cm]
 		 \fint_{\Omega_{35r}}G(|D\psi_i|)\,dz
 		 &\leq \frac{2^{n+1}\lambda\delta}{4}, \quad i=1,2.
 		\end{aligned}
 		\right.
 	\end{equation}
 	To proceed, we introduce the following standard rescaling argument:
 	\begin{equation*} 
 		\left\{
 		\begin{aligned}
 			&\tilde{u}(x) = \frac{u(7rx)}{2^{n+1} \lambda 7r}, \quad \tilde{F}(x) = \frac{F( 7rx)}{2^{n+1} \lambda r},  \quad \tilde{\psi_i}(x) = \frac{\psi_i(7 rx)}{2^{n+1} \lambda 7r}, \quad  i=1,2, \\[0.15cm]
 			&\tilde{a}(x,\eta) = a(7 rx,2^{n+1} \lambda \eta), \quad 	\tilde{b}(x,\eta) = b( 7rx,2^{n+1}\lambda \eta), \quad \tilde{h}(x) = \frac{h(7 rx)}{2^{n+1} \lambda 7r},   \\[0.15cm]
 			&\tilde{g}(t) = g(2^{n+1}\lambda t), \quad \tilde{G}(t) =\int_0^t \tilde{g}(s) ds, \quad t \geq 0.  \\[0.15cm]
 		\end{aligned}
 		\right.
 	\end{equation*} 
 	Under this scaling, conditions \eqref{ass1} and \eqref{ass2} remain satisfied, and we clearly have $\frac{R}{7r} > 5$.
 	Proceeding analogously to the interior case, we apply Lemmas \ref{uw1}, \ref{w1w2l}, \ref{w2w3l}, and \ref{w3w4b}. Scaling back the variables, we deduce the existence of a reference function $w_4 \in W^{1,G}(\Omega_{28r})$ satisfying
 	\begin{align*}
 		\fint_{\Omega_{21r}}
 		|V(Du)-V(Dw_4)|^2\,dz
 		 \leq 2^{n+1}\lambda \varepsilon_2 \quad \& \quad  \|G(|Dw_4|)\|_{L^\infty(\Omega_{14r})}
 		 \leq
 		2^{n+1}\lambda C_2 ,
 	\end{align*}
 where $\varepsilon_2$ is a small parameter to be determined later, and $C_2 > 1$ is a constant depending only on $n, l, L, i_g, s_g$.
 By further enforcing the condition
 	\[
 C\geq \max\{2^{n+2}c_vC_2,27^n\},
 	\]
 	we can argue exactly as in the case $B_{5r}(y) \subset \Omega$ to deduce the set inclusion
 	\begin{align*}
 		\left\{
 		z\in\Omega_{11r}:
 		\mathcal{M}(G(|Du|))>\lambda C
 		\right\}
 		\subset
 		\left\{
 		z\in\Omega_{11r}:
 		\mathcal{M}_{\Omega_{21r}}
 		\left(
 		|V(Du)-V(Dw_4)|^2
 		\right)
 		>2^{n+1}\lambda
 		\right\},
 	\end{align*}
This inclusion naturally yields the measure estimate
 	\begin{align*}
 		\left|
 		\left\{
 		z\in\Omega_{11r}:
 		\mathcal{M}(G(|Du|))>\lambda C
 		\right\}
 		\right|
 		<
 		c\varepsilon_2|B_r(y)|.
 	\end{align*}
 	Finally, observing that $\Omega_r(y) \subset \Omega_{11r}$, we arrive at the desired local bound
 	\begin{align*}
 		\left|
 		\left\{
 		z\in\Omega_r(y):
 		\mathcal{M}(G(|Du|))>\lambda C
 		\right\}
 		\right|
 		<
 		c\varepsilon_2|B_r(y)|.
 	\end{align*}
 The proof is then completed by choosing $\varepsilon_2$ sufficiently small such that $c\varepsilon_2 \leq \varepsilon$.
 \end{proof}

\end{document}
